\section{Related work and positioning}
\label{sec:related}

\subsection{The closest published antecedent}

\textcite{norton2015} is the closest published antecedent to this methodology.
It sets out a framework for prioritising and selecting urban green
infrastructure to mitigate high temperatures, built on a review of the
relationships between urban geometry, green infrastructure, and temperature. Its
logic --- assess the thermal condition of a place, consider who is exposed,
then determine which intervention is suited to that place --- is the logic this
tool implements, and the three-layer structure of \Cref{sec:framework} follows
it directly.

The verification status of this source should be stated plainly, because the
paper leans on it for framing. The bibliographic record was confirmed from the
Monash University research portal, and the framework description was confirmed
from an authoritative secondary description rather than from the full text.
The tool takes no numeric value from \textcite{norton2015}; it takes a
structural argument. \Cref{app:verification} records this alongside every other
source.

Two differences from that antecedent are worth stating. First, this tool is an
operational instrument with a fixed, versioned, machine-readable
parameterisation rather than a conceptual framework, which means every
judgment it embodies must be written down as a number and can therefore be
disputed as a number. Second, it treats climate as a first-order term rather
than as context, for the empirical reason given in \Cref{sec:climate-first}.

\subsection{Composite-indicator practice}

The tool is, structurally, a composite indicator: a set of normalised
sub-indicators combined by weighted aggregation into a headline score. The
construction sequence adopted here --- theoretical framework, variable
selection, treatment of missing data, normalisation, weighting and
aggregation, and presentation --- follows \textcite{nardo2008}, the standard
reference for this class of instrument.

Two of that handbook's positions are load-bearing for this paper and are
adopted without qualification.

The first is that \textbf{weighting is a value choice, not a technical
derivation}. No literature can supply the correct relative importance of
cooling performance against equity against cost, because that ratio expresses a
policy preference rather than a fact about the world. The aggregation weights
in \Cref{sec:aggregation} are therefore declared as expert judgment, versioned,
adjustable in configuration, and are to be defended empirically through the
sensitivity analysis specified in \Cref{sec:sensitivity} rather than by appeal
to authority. Presenting them as derived would be the more persuasive move and
the dishonest one.

The second is that \textbf{uncertainty and sensitivity analysis are integral to
a credible composite indicator, not optional additions}. Because the
aggregation weights are the tool's most consequential unsourced choice, their
influence on the output must be quantified and published rather than asserted
to be small. \Cref{sec:sensitivity} specifies that analysis in full and
\Cref{sec:sensitivity-results} reports its results at version \methver{}.

\subsection{The risk framing}

The Baseline layer follows the standard climate-risk decomposition, in which
risk arises from the interaction of a \emph{hazard} with \emph{exposure} and
\emph{vulnerability}. This is why a physically hot site with no exposed
population scores differently from an equally hot site surrounded by vulnerable
residents, and it is the reason the tool computes a Heat Priority Index rather
than reporting heat alone.

The decomposition is applied concretely: the Heat Exposure Score
(\Cref{sec:heat-exposure}) represents hazard and physical exposure, drawing on
surface-energy-balance drivers for which \textcite{ziter2019} supplies direct
empirical support; the Vulnerability Score (\Cref{sec:vulnerability}) represents
population susceptibility, with its indicator selection grounded in
\textcite{reid2009}; and the two are combined into a single priority index
(\Cref{sec:hpi}).

\subsection{Climate as a first-order term}
\label{sec:climate-first}

The tool's most consequential departure from common screening practice is its
refusal to publish a single global cooling number per intervention type. The
justification is empirical. \textcite{keravec2026} reports that the type of
solution alone explains \SI{12}{\percent} of the variance in observed cooling
effectiveness; adding climate zone raises this to \SI{30}{\percent}; and adding
the K\"oppen--Geiger climate subzone raises it to \SI{78}{\percent}.

\Cref{fig:climate-dependence} shows what that variance decomposition looks like
at the level of individual measures. For green walls, the same synthesis that
reports a \qty{3.0}{\degreeCelsius} global central estimate reports
subzone-specific estimates ranging from \qty{5.7(10)}{\degreeCelsius} in Cwa
and \qty{5.5(10)}{\degreeCelsius} in BWh down to \qty{1.3}{\degreeCelsius} in
Cfa/Cfb and \qty{1.0(8)}{\degreeCelsius} in Af --- an approximately fivefold
spread within a single measure. For street trees, the same review reports
\qty{4.0(16)}{\degreeCelsius} in tropical (A) climates against
\qty{1.7(14)}{\degreeCelsius} in arid (B) climates.

A tool reporting one number per measure would discard this. Climate suitability
is consequently a first-order adjustment in the tool's Layer 2
(\Cref{sec:adjustment}), not a refinement applied afterwards.

\begin{figure}[htbp]
\centering
% Climate dependence of measured cooling effectiveness (keravec2026).
% Left panel: green walls by Koppen-Geiger subzone.
% Right panel: street trees, tropical (A) against arid (B).
% Error bars show the uncertainty the source states; no bar is drawn for
% Cfa/Cfb, for which the source states a central estimate only.
% All values are daytime pedestrian-level AIR temperature reductions.
\begin{tikzpicture}[font=\footnotesize]

\pgfplotsset{
  climatepanel/.style={
    ybar,
    ymin=0, ymax=7.6,
    ytick={0,1,...,7},
    height=64mm,
    axis lines=left,
    axis line style={draw=rulegrey, line width=0.5pt},
    tick style={draw=rulegrey},
    ymajorgrids=true,
    grid style={draw=rulegrey!45, line width=0.3pt, dash pattern=on 1pt off 1.5pt},
    xtick=data,
    xticklabel style={font=\scriptsize},
    error bars/y dir=both,
    error bars/y explicit,
    error bars/error bar style={draw=criterraink, line width=0.7pt},
    nodes near coords,
    nodes near coords style={font=\scriptsize\color{criterraink},
      /pgf/number format/fixed, /pgf/number format/precision=1,
      /pgf/number format/zerofill, yshift=1pt},
    every node near coord/.append style={anchor=south west, xshift=2pt},
  },
}

\begin{axis}[
  name=left,
  climatepanel,
  width=0.56\textwidth,
  bar width=10pt,
  enlarge x limits=0.14,
  ylabel={Air temperature reduction (\si{\degreeCelsius})},
  ylabel style={font=\footnotesize},
  symbolic x coords={Cwa,BWh,Csa,{Cfa/Cfb},Af},
  title={\footnotesize Green walls, by K\"oppen--Geiger subzone},
]
\addplot[draw=criterragreen, fill=criterragreen!75]
  coordinates {
    (Cwa,        5.7) +- (0,1.0)
    (BWh,        5.5) +- (0,1.0)
    (Csa,        2.9) +- (0,0.8)
    ({Cfa/Cfb},  1.3) +- (0,0.0)
    (Af,         1.0) +- (0,0.8)
  };
\draw[dashed, draw=criterraink, line width=0.7pt]
  ({rel axis cs:0,0}|-{axis cs:Cwa,3.0}) -- ({rel axis cs:1,0}|-{axis cs:Cwa,3.0});
\node[anchor=south east, font=\scriptsize\itshape, align=right,
      color=criterraink, inner sep=1.5pt]
  at (axis cs:Af,3.06) {global central estimate\\\SI{3.0}{\degreeCelsius}};
\end{axis}

\begin{axis}[
  name=right,
  climatepanel,
  at={($(left.south east)+(14mm,0)$)}, anchor=south west,
  width=0.40\textwidth,
  bar width=14pt,
  enlarge x limits=0.34,
  yticklabel=\empty,
  symbolic x coords={{Tropical (A)},{Arid (B)}},
  title={\footnotesize Street trees, by climate zone},
]
\addplot[draw=criterragreen, fill=criterragreen!75]
  coordinates {
    ({Tropical (A)}, 4.0) +- (0,1.6)
    ({Arid (B)},     1.7) +- (0,1.4)
  };
\draw[dashed, draw=criterraink, line width=0.7pt]
  ({rel axis cs:0,0}|-{axis cs:{Tropical (A)},3.0})
  -- ({rel axis cs:1,0}|-{axis cs:{Tropical (A)},3.0});
\node[anchor=south east, font=\scriptsize\itshape, align=right,
      color=criterraink, inner sep=1.5pt]
  at (axis cs:{Arid (B)},3.06) {adopted ceiling\\\SI{3.0}{\degreeCelsius}};
\end{axis}

\end{tikzpicture}

\caption[Climate dependence of cooling effectiveness]{Climate dependence of
reported cooling effectiveness, from the umbrella review of
\textcite{keravec2026}. All values are daytime pedestrian-level
\emph{air} temperature reductions. Error bars show the uncertainty the source
states; no bar is drawn for Cfa/Cfb, which this methodology's evidence tables
record as a single combined central estimate without a stated interval. The dashed line marks \qty{3.0}{\degreeCelsius}: the
global central estimate for green walls in the left panel, and the ceiling this
methodology adopts for any single site-level intervention in the right panel.
The spread within each measure is the empirical reason climate suitability is
treated as a first-order adjustment rather than as context.}
\label{fig:climate-dependence}
\end{figure}

\subsection{By what standard should a screening tool be judged?}
\label{sec:standards}

A recurring failure in the reception of screening instruments is that they are
evaluated against simulation standards, which they cannot meet and were not
built to meet. If the question is ``does this reproduce ENVI-met to within
\qty{0.2}{\degreeCelsius}?'', the answer is no, and it would be no for any
instrument that asks forty-five questions and runs in seconds. That is not a
finding about this tool; it is a restatement of what a screening instrument is.

The productive questions are different, and this paper invites reviewers to
apply them.

\begin{enumerate}
  \item \textbf{Is every value traceable?} Can a reader locate the source of
  every number the tool applies, and see what was actually verified about that
  source? \Cref{app:typologies,app:verification} exist to make this checkable.

  \item \textbf{Are the uncertainties honestly represented?} Does the tool
  report ranges that reflect the spread in the underlying evidence, and does it
  lower its stated confidence when it knows less? \Cref{sec:uncertainty}
  specifies the mechanism.

  \item \textbf{Is the ranking stable?} Because prioritisation is the tool's
  actual function, the relevant question about the weights is not whether they
  are correct --- they are a value choice --- but whether the ordering they
  produce survives reasonable perturbation. \Cref{sec:sensitivity} specifies
  how this will be measured and published.

  \item \textbf{Are the value judgments disclosed?} Where the methodology takes
  a position that a reasonable person could reject --- most notably the
  equity-forward weighting of \Cref{sec:effective-weights} --- is that position
  stated openly and quantified, or is it buried in an aggregation?

  \item \textbf{Does it refuse to answer when it should?} An instrument that
  always produces a number is less trustworthy than one that reports
  \emph{not estimated} when its inputs do not support an estimate. The cost
  policy of \Cref{sec:costs} and the missing-data rules of
  \Cref{sec:missing-data} are deliberate refusals.
\end{enumerate}

A screening tool that satisfies these five is doing its job. One that fails
them is not rescued by agreeing with a simulation on a test case.
